\item (a) Utilice coordenadas cilíndricas para demostrar que el volumen del sólido limitado por arriba por la esfera $r^2 + z^2 = a^2$ y por abajo por el cono $z = r \cot \phi_0$ (o $\phi = \phi_0$), en donde $0 < \phi_0 < \pi/2$, es
\[ V = \frac{2\pi a^3}{3}(1 - \cos \phi_0) \]
(b) Deduzca que el volumen de la cuña esférica descrita por $\rho_1 \le \rho \le \rho_2, \ \theta_1 \le \theta \le \theta_2, \ \phi_1 \le \phi \le \phi_2$ es
\[ \Delta V = \frac{\rho_2^3 - \rho_1^3}{3} (\cos \phi_1 - \cos \phi_2)(\theta_2 - \theta_1) \]
(c) Aplique el teorema del valor medio para demostrar que el volumen considerado en la parte (b) se puede escribir como
\[ \Delta V = \tilde{\rho}^2 \operatorname{sen} \tilde{\phi} \ \Delta\rho \ \Delta\theta \ \Delta\phi \]
en donde $\tilde{\rho}$ se encuentra entre $\rho_1$ y $\rho_2$, $\tilde{\phi}$ se encuentra entre $\phi_1$ y $\phi_2$, $\Delta\rho = \rho_2 - \rho_1$, $\Delta\theta = \theta_2 - \theta_1$, y $\Delta\phi = \phi_2 - \phi_1$.
